% Generated draft: requires mathematical and source review.
\documentclass[11pt]{article}
\usepackage[margin=25mm]{geometry}
\usepackage{fontspec}
\setmainfont{Latin Modern Roman}
\usepackage{amsmath,amssymb,mathtools,unicode-math}
\setmathfont{Latin Modern Math}
\usepackage{enumitem,needspace,xurl,hyperref}
\setlength{\parindent}{0pt}
\setlength{\parskip}{5pt}
\setlength{\emergencystretch}{3em}
\newcommand{\N}{\mathbb{N}}
\newcommand{\Z}{\mathbb{Z}}
\newcommand{\Q}{\mathbb{Q}}
\newcommand{\R}{\mathbb{R}}
\newcommand{\C}{\mathbb{C}}
\newcommand{\F}{\mathbb{F}}
\newcommand{\A}{\mathbb{A}}
\newcommand{\PP}{\mathbb{P}}
\newcommand{\E}{\mathbb{E}}
\newcommand{\eps}{\varepsilon}
\newcommand{\dd}{\,\mathrm{d}}
\begin{document}
\section*{UnsolvedMath ID: 800367}
\textbf{Source problem code:} catalogue entry
\subsection{Definitions and mathematical statement}
Let \((W,S)\) be a Coxeter system. The supplied record does not define the relation \(\leq\). Under the conventional interpretation as strong Bruhat order, \(u\leq v\) means that a reduced expression for \(v\) contains a subword that is a reduced expression for \(u\).

Let \(c,d\) be noncommuting indeterminates, and let \(\{c,d\}^{*}\) denote the set of finite words in them, including the empty word. Define
\[
\mathbb{Z}_{\geq 0}\langle c,d\rangle
=
\left\{\sum_{w\in\{c,d\}^{*}} a_w w:
 a_w\in\mathbb{Z}_{\geq 0},\quad
 a_w=0\text{ for all but finitely many }w\right\},
\]
with multiplication induced by concatenation of words. This makes explicit the nonnegative-coefficient interpretation of the source notation \(\mathbb{N}\langle c,d\rangle\).

\textbf{Conjecture 5 (Brenti, 2024; conditional transcription).}
For every Coxeter system \((W,S)\) and every \(u,v\in W\) with \(u\leq v\),
\[
\widetilde{\Phi}_{u,v}\in\mathbb{Z}_{\geq 0}\langle c,d\rangle.
\]
Thus, if the source invariant is written as a noncommutative polynomial
\[
\widetilde{\Phi}_{u,v}=\sum_{w\in\{c,d\}^{*}} a_w(u,v)w,
\]
the assertion is that every coefficient \(a_w(u,v)\) is a nonnegative integer.

The supplied record does not define \(\widetilde{\Phi}_{u,v}\), its construction, or its normalization. Consequently, it does not by itself determine the invariant whose coefficients are asserted to be nonnegative. The Bruhat-order interpretation above also requires confirmation from the cited source. No finiteness restriction on \(W\) or \(S\) is stated. The quantifier includes \(u=v\), but the value or convention for \(\widetilde{\Phi}_{u,u}\) is not supplied.
\subsection{Short English statement}
Brenti's supplied Conjecture 5 asserts that, for every comparable pair in every Coxeter system, \(\widetilde{\Phi}_{u,v}\) is a polynomial in noncommuting \(c,d\) with nonnegative integer coefficients. The invariant is undefined in the supplied record, and interpreting comparability as strong Bruhat order requires confirmation. This is a statement-only draft requiring human mathematical review.
\subsection{Sources}
\begin{itemize}
\item Supplied local UnsolvedMath registry record: permanent ID \(800367\), source problem code \texttt{catalogue entry}, title ``Coxeter Groups and Unimodality --- Conjecture 5,'' in the AMR Open Problem Lists collection. The record attributes the conjecture to Brenti and gives the proposed year as 2024.
\item Brenti, \emph{Some Open Problems on Coxeter Groups and Unimodality} (2024), Conjecture 5. Supplied source URL: \url{https://arxiv.org/abs/2410.09897}. The record reports access on 2026-10-16 and extraction by \texttt{arxiv\_latex}; these are supplied provenance statements, not independently verified access or extraction.
\item The local record labels the problem ``open'' and records an update timestamp of \texttt{2026-10-16T14:00:00Z}. This is supplied registry curation, not an independently verified literature assessment.
\end{itemize}
\end{document}
